% wave13_a15_mathieu_conj1_segal.tex
% Wave 13, A15. Adversarial assault and rigorous healing of the
% Mathieu-moonshine / Conjecture 1 side of Lorgat 2020.
%
% Voice: Graeme Segal (elliptic cohomology, modular functors, CFT
% axioms) + Greg Moore (BPS state algebra, Moore--Tachikawa precision,
% automorphic / arithmetic discipline on BKM).
%
% Primary sources, read only through
%   /Users/raeez/calabi-yau-quantum-groups/notes/wave13_source_paper.txt
%   /Users/raeez/calabi-yau-quantum-groups/notes/wave13_source_slides.txt
% (the user's 2020 paper and its companion slide deck).
%
% Cross-reference:
%   wave12_a15_moonshine_segal.tex (moonshine / Niemeier atlas);
%   wave12_a5_mock_modular_polyakov.tex (mock modular K3 at d=2,
%     Gannon positivity, 25-vs-21 stratification).
%
% Complementary: wave9_p5_mathieu_26_class_table.tex (canonical
% ATLAS numerics); wave10_r2_mathieu_1_g_twined.tex (four-step
% waterfall at (1,g)); wave9_m4_clery_gritsenko_fourier_check.tex
% (paramodular Fourier audit).
%
% Format: read-only vs working_notes.tex and the manuscript;
% write-only this file. Internal working-notes artefact.
% Not part of the manuscript. Raeez Lorgat author.

\documentclass[11pt]{article}
\usepackage[margin=1in]{geometry}
\usepackage{amsmath, amssymb, amsthm, mathtools, enumitem, booktabs, longtable}
\usepackage{hyperref}

\newtheorem{theorem}{Theorem}
\newtheorem{proposition}[theorem]{Proposition}
\newtheorem{lemma}[theorem]{Lemma}
\newtheorem{corollary}[theorem]{Corollary}
\newtheorem{conjecture}[theorem]{Conjecture}
\newtheorem{definition}[theorem]{Definition}
\newtheorem{observation}[theorem]{Observation}
\newtheorem{warning}[theorem]{Warning}
\newtheorem*{attack}{Attack}
\newtheorem*{heal}{Heal}
\newtheorem*{verification}{Verification}

\providecommand{\ClaimStatusTheorem}{\textbf{[Theorem]}\ }
\providecommand{\ClaimStatusConjectured}{\textbf{[Conjectured]}\ }
\providecommand{\ClaimStatusHeuristic}{\textbf{[Heuristic]}\ }
\providecommand{\ClaimStatusConditional}[1]{\textbf{[Conditional on: #1]}\ }
\providecommand{\ClaimStatusFalse}{\textbf{[False]}\ }

\providecommand{\kch}{\kappa_{\mathrm{ch}}}
\providecommand{\kcat}{\kappa_{\mathrm{cat}}}
\providecommand{\kBKM}{\kappa_{\mathrm{BKM}}}
\providecommand{\kfib}{\kappa_{\mathrm{fiber}}}

\providecommand{\ZZ}{\mathbb{Z}}
\providecommand{\QQ}{\mathbb{Q}}
\providecommand{\CC}{\mathbb{C}}
\providecommand{\RR}{\mathbb{R}}
\providecommand{\NN}{\mathbb{N}}
\providecommand{\HH}{\mathbb{H}}
\providecommand{\cO}{\mathcal{O}}
\providecommand{\cZ}{\mathcal{Z}}
\providecommand{\cH}{\mathcal{H}}
\providecommand{\fg}{\mathfrak{g}}
\providecommand{\fm}{\mathfrak{m}}

\providecommand{\Sp}{\mathrm{Sp}}
\providecommand{\SL}{\mathrm{SL}}
\providecommand{\SO}{\mathrm{SO}}
\providecommand{\GL}{\mathrm{GL}}
\providecommand{\Aut}{\mathrm{Aut}}
\providecommand{\Sym}{\mathrm{Sym}}
\providecommand{\tr}{\mathrm{tr}}
\providecommand{\wt}{\mathrm{wt}}
\providecommand{\Tot}{\mathrm{Tot}}
\providecommand{\Frame}{\mathrm{Frame}}
\providecommand{\mult}{\mathrm{mult}}
\providecommand{\DT}{\mathrm{DT}}
\providecommand{\red}{\mathrm{red}}

\providecommand{\Delf}{\Delta_5}
\providecommand{\phipq}{\phi_{0,1}}
\providecommand{\Mtf}{M_{24}}
\providecommand{\Zg}{Z^{(g)}_{\mathrm{K3}}}
\providecommand{\Lamtt}{\Lambda^{2,1}_{I\!I}}

\title{The Mathieu side of Conjecture~1:\\
  adversarial assault and rigorous healing\\[0.3em]
  \textit{(Segal voice: modular functors, elliptic-cohomology
  $\pi_0 \mathrm{Ell}(X) = $ weak Jacobi forms; Moore voice: BPS
  state algebra, automorphic precision, GKM denominator identity)}}
\author{Wave 13, note A15 (internal; \texttt{notes/} only)}
\date{2026-04-22}

\begin{document}
\maketitle

\begin{abstract}
We subject the Mathieu-moonshine / Conjecture~1 side of Lorgat 2020
to eight adversarial attack--heal cycles. The target object is the
alleged ``correspondence'' between the \emph{8 Gritsenko--Cl\'ery
diagonal-divisor paramodular forms} (Gritsenko--Cl\'ery
\emph{J.\ Reine Angew.\ Math.}\ 2011, arXiv:0812.3962, Thm.\ 1.1) and
the \emph{26 twined K3 elliptic genera} for $g \in \Mtf$
(Eguchi--Ooguri--Tachikawa 2011 \emph{Exp.\ Math.}\ 20;
Gaberdiel--Hohenegger--Volpato 2012 arXiv:1211.7074), as advertised in
Lorgat 2020 \S6 Conjecture~1. The adversarial line is that three
\emph{numerically distinct} sets compete for the role of
``admissible~$N$'' --- namely $\{1, 2, 3, 4\}$ (the Gritsenko--Cl\'ery
paramodular levels, supplying eight forms across $t \in \{1, 2, 3, 4\}$),
$\{1, 2, 3, 4, 6\}$ (the CHL / genus-two Borcherds admissible levels of
Gritsenko--Cl\'ery 2013 + Oberdieck 2018 Igusa cusp conjecture), and
$\{1, 2, 3, 4, 5, 6, 7, 8\}$ (the Nikulin-Mukai symplectic-automorphism
range for K3) --- and Lorgat 2020 \S6 Conjecture~1 does \emph{not}
clearly specify which. The paper asserts ``all eight diagonal-divisor
modular forms'' arise as reciprocal-square-roots of twisted DT zeta
functions indexed by $(g_N, h_M)$ with $N, M \le 8$, which would
suggest the Nikulin-Mukai range $1 \le N \le 8$ --- but only $N \in
\{1, 2, 3, 4, 6\}$ have positive-weight Gritsenko lifts, only $N = 1$
is proven (Oberdieck--Pixton--Pandharipande 2016--2018), and only 4
of the 8 Gritsenko--Cl\'ery forms sit over CHL-admissible $N$; the
remaining 4 sit over $(t, N) \in \{(2,1), (1,3), (1,4), (4,1),
(1,2{\cdot}2)\}$ combinations whose BKM / elliptic-genus image is not
directly computed. We heal by stratifying Conjecture~1 into five
rigorously scoped sub-conjectures, pinning each Gritsenko--Cl\'ery
form to its $(t, N, \nu)$ datum and its admissible $M_{24}$-class
image, and flagging which layer is theorem, which conjecture,
which speculative. Net: Conjecture~1 is a theorem at
$\{(t, N)\} = \{(1, 1)\}$ ($\Delf$ / class~$1A$) and a
well-posed arithmetic correspondence at $\{(2, 1), (1, 2), (1, 3),
(1, 4)\}$; the remaining 3 GC forms and the 16 Mukai-unlifted
$M_{24}$-classes remain genuinely open, with precise obstructions
named at each stratum.
\end{abstract}

\tableofcontents

% ============================================================
\section{Three competing admissibility sets, one Conjecture~1}
\label{sec:three-sets}
% ============================================================

\subsection{The three sets}

The Mathieu side of Lorgat 2020 Conjecture~1 (paper \S6,
\emph{Arithmetic Geometry of some CY$_3$'s}) is articulated against
three genuinely different admissibility sets that the paper does not
distinguish. Segal's modular-functor axiomatisation forces the
distinction: different admissibility sets label different modular
functors on different cobordism categories.

\begin{enumerate}[label=(\Alph*)]
\item\textbf{The Gritsenko--Cl\'ery paramodular set} $\mathsf{GC}$.
The Gritsenko--Cl\'ery 2011 theorem (our paper Thm.~1, following
arXiv:0812.3962) enumerates Siegel modular forms on $\Gamma_t(N)$
vanishing on the translates of the diagonal to first order. The total
count is 8, distributed across paramodular levels $(t, N)$ as

\begin{center}\small
\begin{tabular}{c|c|c|c}
\toprule
Form (paper labelling) & $\Gamma_t(N)$ & Weight & Character \\
\midrule
$\Delf$ (Igusa) & $\Gamma_1 = \Sp_4(\ZZ)$ & $5$ & $\nu^2$ \\
$F_{2, \Gamma_2}$ & $\Gamma_2$ & $2$ & $\nu^4$ \\
$F_{3, \Gamma_1(2)}$ & $\Gamma_1(2)$ & $3$ & $\nu^2$ \\
$F_{1, \Gamma_3}$ & $\Gamma_3$ & $1$ & $\nu^6$ \\
$F_{2, \Gamma_1(3)}$ & $\Gamma_1(3)$ & $2$ & $\nu^2$ \\
$F_{1/2, \Gamma_4}$ & $\Gamma_4$ & $1/2$ & $\nu^8$ \\
$F_{3/2, \Gamma_1(4)}$ & $\Gamma_1(4)$ & $3/2$ & $\nu^4$ \\
$F_{1, \Gamma_2(2)}$ & $\Gamma_2(2)$ & $1$ & $\nu^4$ \\
\bottomrule
\end{tabular}
\end{center}

Eight forms, distributed over paramodular-index parameters
$t \in \{1, 2, 3, 4\}$ and Hecke-type congruence parameters $N \in
\{1, 2, 3, 4\}$. The pair $(t, N)$ is the fundamental datum; the
``order~$N$''-counting fiction conflates $t$ and $N$.

\item\textbf{The CHL / Borcherds-admissible set} $\mathsf{CHL} = \{1,
2, 3, 4, 6\}$. Gritsenko--Cl\'ery 2013 \emph{Compos.\ Math.}\ 149
(arXiv:1302.0272) classifies the Siegel paramodular cusp forms whose
Borcherds-lift expansion produces a rank-$3$ BKM at a standard
prime-level CHL orbifold of the K3 SCFT. The five admissible orders
are $1, 2, 3, 4, 6$, with Borcherds weights $\kBKM = c_N(0)/2 = (5, 4,
3, 2, 1)$ respectively.

\item\textbf{The Nikulin-Mukai range} $\mathsf{NM} = \{1, 2, 3, 4, 5,
6, 7, 8\}$. Nikulin 1979--1984 (\emph{Proc.\ Steklov Inst.\ Math.}\
165) and Mukai 1988 (\emph{Invent.\ Math.}\ 94) classify finite
symplectic automorphism groups of K3, enumerating $N \le 8$ for
cyclic orders. Lorgat 2020 \S6 explicitly mentions the bound
$N, M \le 8$ ``by a theorem of Nikulin''.
\end{enumerate}

\subsection{The three sets do not agree}

These three sets have cardinalities $8, 5, 8$ respectively (as
multisets), and they index three \emph{different} arithmetic
classifications:

\[
\underbrace{|\mathsf{GC}| = 8}_{\text{paramodular forms}} \;\ne\;
\underbrace{|\mathsf{CHL}| = 5}_{\text{Borcherds-admissible orders}}
\;\ne\; \underbrace{|\mathsf{NM}| = 8}_{\text{Nikulin-Mukai auto orders}}.
\]

(Numerically $|\mathsf{GC}| = |\mathsf{NM}| = 8$; \emph{the equality
is a coincidence}, not a canonical bijection. $\mathsf{GC}$ is
indexed by $(t, N)$ with $t \in \{1,2,3,4\}, N \in \{1,2,3,4\}$
subject to diagonal-divisor discriminant; $\mathsf{NM}$ is indexed by
cyclic orders $1, 2, \ldots, 8$.)

The intersection structure is:

\begin{itemize}
\item $\mathsf{GC} \cap \mathsf{CHL}$ (forms whose paramodular level
agrees with a CHL-admissible order): the four forms with
$N \in \{1, 2, 3, 4\}$ on $\Gamma_1(N)$ or $\Gamma_t$ --- namely
$\Delf$, $F_{2, \Gamma_2}$, $F_{3, \Gamma_1(2)}$, $F_{1, \Gamma_3}$,
$F_{2, \Gamma_1(3)}$, $F_{3/2, \Gamma_1(4)}$ --- with precise
matching requiring specification of the level.
\item The $N = 6$ CHL case has \emph{no} Gritsenko--Cl\'ery
diagonal-divisor form on $\Gamma_t(6)$ (the $N = 6$ Borcherds lift at
weight 1 comes from the Gritsenko 1999 construction, not from a
diagonal-divisor form).
\item The four $\mathsf{GC}$ forms with parameters not in $\mathsf{CHL}$
are unaccounted for in the Borcherds-BKM dictionary of CHL.
\end{itemize}

\subsection{Conjecture~1 in its precise category: the map that Lorgat
2020 \S6 is actually asserting}

Lorgat 2020 \S6 p.~2 Conjecture~1 makes two assertions. First, all
eight diagonal-divisor Gritsenko--Cl\'ery paramodular forms arise,
up to a constant $C$, as reciprocal-square roots of the twisted DT
zeta function $Z^X_{L, h_M}$. Second, each of these Siegel paramodular
forms is the denominator function of a generalised
Borcherds--Kac--Moody superalgebra whose root multiplicities are
specified by the $(g_N, h_M)$-twisted twined elliptic genus of a K3
surface.

The conjecture is a four-layer statement, one layer per bullet:

\begin{enumerate}[label=(L\arabic*)]
\item (DT zeta / Siegel layer) each of the 8 Gritsenko--Cl\'ery forms
$\Phi \in \mathsf{GC}$ satisfies
\[
Z^X_{L, h_M}(q, t, p) \;=\; -\, C(\Phi) \,/\, \Phi^2
\]
for an explicit CY$_3$ $X = (S \times E)/(\ZZ/N\ZZ)$ with
$(S, L, g_N, h_M)$ data and an explicit constant $C(\Phi)$.
\item (BKM denominator layer) each $\Phi \in \mathsf{GC}$ is the
denominator of a rank-$3$ generalised BKM superalgebra $\fg_\Phi$.
\item (Mathieu layer) the root multiplicities $\mult_\Phi(\alpha)$ are
equal to the Fourier coefficients of the $(g_N, h_M)$-twisted twined
elliptic genus of K3.
\item (Enumeration closure) this accounts for \emph{all 8} of
$\mathsf{GC}$, with no additional structure beyond $(N, M)$.
\end{enumerate}

Segal's axiomatisation demands, at each layer, a functor landing in a
precise target category: $\mathsf{ParMod}(\Gamma_t(N)) \to$ (DT
generating series on K3$\times$E moduli) at (L1);
$\mathsf{ParMod}(\Gamma_t(N)) \to \mathsf{BKM}^{(3)}$ at (L2);
$\mathsf{BKM}^{(3)} \to \mathsf{JacForms}_{0, 1}(\Gamma_0(N_g))$ at
(L3). Each of these functors must be stated separately.

% ============================================================
\section{Primary data: the 26 classes, the 8 forms, the intended map}
\label{sec:primary-data}
% ============================================================

\subsection{ATLAS 26 classes of $\Mtf$ with Frame shape}

The Mathieu group $\Mtf$ has exactly 26 conjugacy classes (not 25;
see the cross-reference warning in \S~\ref{subsec:25-vs-26} below).
In ATLAS order:

\begin{center}\footnotesize
\begin{tabular}{l|l|c|c|c|c|c|c}
\toprule
$[g]$ & cycle shape & $|g|$ & $\chi_{24}(g)$ & $c_N(0)$ & Frame $\eta_g$ &
  Mukai-realisable? & GC form? \\
\midrule
$1A$   & $1^{24}$              & 1  & 24 & 10 & $\eta(\tau)^{24}$ & yes & $\Delf$ \\
$2A$   & $1^8 2^8$             & 2  & 8  &  8 & $\eta(\tau)^8 \eta(2\tau)^8$ & yes & $F_{2,\Gamma_2}$ \\
$2B$   & $2^{12}$              & 2  & 0  & $-4$& $\eta(2\tau)^{12}$ & yes & none (wt $\le 0$) \\
$3A$   & $1^6 3^6$             & 3  & 6  &  6 & $\eta(\tau)^6 \eta(3\tau)^6$ & yes & $F_{1,\Gamma_3}$ \\
$3B$   & $3^8$                 & 3  & 0  & $-4$& $\eta(3\tau)^8$ & yes & none \\
$4A$   & $2^4 4^4$             & 4  & 0  &  0 & $\eta(2\tau)^4 \eta(4\tau)^4$ & yes & none (wt $0$) \\
$4B$   & $1^4 2^2 4^4$         & 4  & 4  &  4 & $\eta(\tau)^4 \eta(2\tau)^2 \eta(4\tau)^4$ & yes & $F_{3/2,\Gamma_1(4)}?$ \\
$4C$   & $4^6$                 & 4  & 0  & $-4$& $\eta(4\tau)^6$ & yes & none \\
$5A$   & $1^4 5^4$             & 5  & 4  &  2 & $\eta(\tau)^4 \eta(5\tau)^4$ & yes & none \\
$6A$   & $1^2 2^2 3^2 6^2$     & 6  & 2  &  2 & $\eta(\tau)^2 \eta(2\tau)^2 \eta(3\tau)^2 \eta(6\tau)^2$ & yes & none \\
$6B$   & $6^4$                 & 6  & 0  &  0 & $\eta(6\tau)^4$ & yes & none \\
$7A$   & $1^3 7^3$             & 7  & 3  & $-1$& $\eta(\tau)^3 \eta(7\tau)^3$ & yes$^\dagger$ & none \\
$7B$   & $1^3 7^3$             & 7  & 3  & $-1$& $\eta(\tau)^3 \eta(7\tau)^3$ & yes$^\dagger$ & none \\
$8A$   & $1^2 2 \cdot 4 \cdot 8^2$ & 8  & 2  &  0 & $\eta(\tau)^2 \eta(2\tau) \eta(4\tau) \eta(8\tau)^2$ & yes & none \\
$10A$  & $2^2 10^2$            & 10 & 0  & $-4$& $\eta(2\tau)^2 \eta(10\tau)^2$ & no & none \\
$11A$  & $1^2 11^2$            & 11 & 2  & $-2$& $\eta(\tau)^2 \eta(11\tau)^2$ & yes$^\dagger$ & none \\
$12A$  & $2\cdot 4\cdot 6\cdot 12$ & 12 & 0 &  0 & $\eta(2\tau)\eta(4\tau)\eta(6\tau)\eta(12\tau)$ & no & none \\
$12B$  & $12^2$                & 12 & 0  &  0 & $\eta(12\tau)^2$ & no & none \\
$14A$  & $1\cdot 2\cdot 7\cdot 14$ & 14 & 2 &  0 & $\eta(\tau)\eta(2\tau)\eta(7\tau)\eta(14\tau)$ & no & none \\
$14B$  & $1\cdot 2\cdot 7\cdot 14$ & 14 & 2 &  0 & $\eta(\tau)\eta(2\tau)\eta(7\tau)\eta(14\tau)$ & no & none \\
$15A$  & $1\cdot 3\cdot 5\cdot 15$ & 15 & 2 &  0 & $\eta(\tau)\eta(3\tau)\eta(5\tau)\eta(15\tau)$ & no & none \\
$15B$  & $1\cdot 3\cdot 5\cdot 15$ & 15 & 2 &  0 & $\eta(\tau)\eta(3\tau)\eta(5\tau)\eta(15\tau)$ & no & none \\
$21A$  & $3\cdot 21$           & 21 & 0  & $-1$& $\eta(3\tau)\eta(21\tau)$ & no & none \\
$21B$  & $3\cdot 21$           & 21 & 0  & $-1$& $\eta(3\tau)\eta(21\tau)$ & no & none \\
$23A$  & $1\cdot 23$           & 23 & 2  &  0 & $\eta(\tau)\eta(23\tau)$ & yes$^\dagger$ & none \\
$23B$  & $1\cdot 23$           & 23 & 2  &  0 & $\eta(\tau)\eta(23\tau)$ & yes$^\dagger$ & none \\
\bottomrule
\end{tabular}
\end{center}

Notation: $\chi_{24}(g)$ is the character of $g$ in the defining
$\mathbf{24}$-representation of $\Mtf$ (equivalently the number of
fixed points in the Golay-code action, equivalently the leading
coefficient of the twined Jacobi form); $c_N(0) := c^{(g)}(0, 0)$ is
the DVV-normalised $q^0 y^0$ constant of the twined weak Jacobi form
$\phi^{(g)}_{0, 1}(\tau, z)$, equal to $\chi_{24}(g)/2$ plus a
correction from the $q^0 y^{\pm 1}$ slots (see Wave~9 P5 Table and
Gaberdiel--Hohenegger--Volpato 2012 Table~4). ``Mukai-realisable''
records whether $g$ occurs as a symplectic automorphism of some K3
surface (Mukai 1988); the five dagger-marked entries $(7A, 7B, 11A,
23A, 23B)$ are realisable only after extension to Mathieu-moonshine
sigma-model CFT symmetry (Gaberdiel--Hohenegger--Volpato 2011
arXiv:1008.3778 \S4). ``GC form?'' names a plausible matching
Gritsenko--Cl\'ery paramodular form (see next subsection for the
matching heuristic).

\subsection{Distinction between 25 and 26: a recurring cross-reference
drift}
\label{subsec:25-vs-26}

\begin{warning}[ATLAS 26, not 25]
\label{warn:26-atlas}
The Mathieu group $\Mtf$ has exactly 26 conjugacy classes (Conway et
al., \emph{ATLAS of Finite Groups}, 1985, p.~96; Thompson 1979).
The 26 are:
\[
1A,\ 2A,\ 2B,\ 3A,\ 3B,\ 4A,\ 4B,\ 4C,\ 5A,\ 6A,\ 6B,\ 7A,\ 7B,\ 8A,
\]
\[
10A,\ 11A,\ 12A,\ 12B,\ 14A,\ 14B,\ 15A,\ 15B,\ 21A,\ 21B,\ 23A,\ 23B.
\]
Some sources (including \texttt{wave12\_a5\_mock\_modular\_polyakov.tex}
\S\ref{sec:cycle3} as well as a handful of secondary references, and the
current Wave-13 prompt) quote \emph{25}, conflating the Galois-paired
classes $\{23A, 23B\}$ into a single abstract ``class
$23A\!/B$''. Galois-paired classes \emph{are distinct as $\Mtf$
conjugacy classes}; they are mapped to one another by a Galois
automorphism of the $23$-cyclotomic field, but each has its own
cycle-shape representative and centraliser. The correct count is 26.
This note uses 26 throughout.
\end{warning}

\subsection{The 8 Gritsenko--Cl\'ery forms and the intended 8-class map}
\label{subsec:8-to-8}

The 8 forms live over levels $\Gamma_t(N) \subset \Sp_4(\ZZ)$ with
$t \in \{1, 2, 3, 4\}$, $N \in \{1, 2, 3, 4\}$, multiplier systems
$\nu^k$ with $k \in \{2, 4, 6, 8\}$. Lorgat 2020 \S6 Conjecture~1
asserts a \emph{correspondence} between the 8 forms and the K3 data
$(g_N, h_M)$ with $N, M \le 8$, mediated by the twisted DT zeta
function and the twisted twined elliptic genus. But which $(g_N, h_M)$
pairs?

The paper provides no explicit 8-element list of $(g_N, h_M)$; it says
only $N, M \le 8$. This leaves multiple possible pairings. The two
natural candidates:

\begin{itemize}
\item[(i)] \textbf{Diagonal pairing} $(g_N, \mathrm{id})$ for
$N \in \{1, 2, 3, 4, 6, 8\}$ plus two off-diagonal pairs: six classes
diagonal ($1A, 2A, 3A, 4B, 6A, 8A$) plus two pairs $(g_2, h_2)$ and
$(g_4, h_4)$. This matches the $\mathsf{NM}$ bound $\le 8$ with the
proviso that only Mukai-realisable $N$ occur. Eight items total.
\item[(ii)] \textbf{Paramodular pairing} one per Gritsenko--Cl\'ery
form, indexed by $(t, N)$; the $(g_{t}, h_{N})$-twisted twined elliptic
genus provides the Fourier multiplicities. Eight items, matching the
$(t, N)$ parameter space above.
\end{itemize}

Candidate (ii) is the more natural reading, since the Gritsenko--Cl\'ery
form \emph{itself} is tagged by $(t, N)$, and the proposed Borcherds lift
at paramodular level $\Gamma_t(N)$ reads off the twined-$g_N$ seed Jacobi
form at fibre orbifold of order $N$ and base orbifold of order $t$.

\textbf{This is the first attack surface.} Lorgat 2020 does not select
between (i) and (ii); the user (at the time of authorship) did not
observe this discrepancy. We heal by picking (ii) and stating
Conjecture~1 with explicit $(t, N, g, h)$-assignments.

% ============================================================
\section{Attack--heal cycle 1: admissible-$N$ ambiguity}
\label{sec:cycle1}
% ============================================================

\subsection*{Scope}

Target: the role of ``$N$'' in Lorgat 2020 Conjecture~1 is triple-overloaded
(paramodular index $t$, Hecke level $N$ of $\Gamma_t(N)$, K3
automorphism order). The three sets of \S\ref{sec:three-sets} are
distinct.

\begin{attack}[Which $N$ is the one in Conjecture~1?]
Lorgat 2020 \S6 uses the symbol $N$ in three distinct senses in two
adjacent paragraphs:
\begin{itemize}
\item The CHL quotient $X = (S \times E)/(\ZZ/N\ZZ)$ uses $N$ as the
K3-automorphism order of $g_N$.
\item The Gritsenko--Cl\'ery Theorem~1 uses $N$ as the Hecke-type
congruence parameter in $\Gamma_t(N)$.
\item The ``$N \le 8$ by Nikulin'' bound uses $N$ as the Nikulin-Mukai
symplectic-automorphism order.
\end{itemize}
These are three different variables. The first ranges over $\{1, 2, 3,
4, 5, 6, 7, 8\}$ (Mukai); the second ranges over $\{1, 2, 3, 4\}$
(Gritsenko--Cl\'ery Theorem~1); the third subsumes the first
($\mathsf{NM}$). The conjecture as stated does not specify which
identification of the 8 forms with 8 data points is intended. As a
result, the conjecture is not well-posed as written.
\end{attack}

\begin{heal}[Stratify into five well-posed sub-conjectures]
We separate the map into five strata, distinguishing paramodular index
$t$, Hecke level $N$ (denoted $N_{\Gamma}$), K3-fibre automorphism
order $N_g$, base-elliptic-curve automorphism order $M_h$, and
$\Mtf$-class $[g]$.

\begin{conjecture}[\ClaimStatusConjectured\ Conjecture~1.1:
paramodular $\leftrightarrow$ BKM]
\label{conj:1-1-param-bkm}
For each of the 8 Gritsenko--Cl\'ery forms $\Phi_{(t, N_\Gamma)} \in
\mathsf{GC}$ with paramodular data $(t, N_\Gamma)$, there exists a
rank-3 generalised Borcherds--Kac--Moody Lie superalgebra
$\fg_{\Phi_{(t, N_\Gamma)}}$ whose denominator function is
$\Phi_{(t, N_\Gamma)}^2$ (or $\Phi$ itself, depending on multiplier
normalisation).
\end{conjecture}

\begin{conjecture}[\ClaimStatusConjectured\ Conjecture~1.2:
BKM $\leftrightarrow$ twined genus]
\label{conj:1-2-bkm-twined}
The root multiplicities of $\fg_{\Phi_{(t, N_\Gamma)}}$ are the Fourier
coefficients of the $(g_{N_g}, h_{M_h})$-twisted twined K3 elliptic
genus, for an explicit pair $(N_g, M_h)$ determined by $(t, N_\Gamma)$.
\end{conjecture}

\begin{conjecture}[\ClaimStatusConjectured\ Conjecture~1.3:
DT side]
\label{conj:1-3-dt}
For each $(t, N_\Gamma) \in \mathsf{GC}$, the reduced DT zeta function
of the CHL quotient $X = (S \times E)/(\ZZ/N_g\ZZ)$ satisfies
\[
Z^X_{\red, L, h_{M_h}}(q, t, p) \;=\; -\, C_{(t, N_\Gamma)} \,/\,
\Phi_{(t, N_\Gamma)}(q, t, p)^2.
\]
\end{conjecture}

\begin{conjecture}[\ClaimStatusConjectured\ Conjecture~1.4:
enumeration closure]
\label{conj:1-4-closure}
The eight pairs $\{(t, N_\Gamma)\}$ of Gritsenko--Cl\'ery parameters
exhaust the parameter space on which Conjectures 1.1--1.3 hold;
equivalently the $(g, h)$-twining pairs of Conjecture~1.2 are exactly
eight, matching $\mathsf{GC}$ in bijection.
\end{conjecture}

\begin{conjecture}[\ClaimStatusConjectured\ Conjecture~1.5:
Mathieu compatibility]
\label{conj:1-5-mathieu}
For $(t, N_\Gamma) = (1, 1)$ and the $(g_{N_g}, h_{M_h})$ pair
determined by Conjecture~1.2, the root multiplicity
$\mult_{\Phi}(\alpha)$ is in fact a non-negative integer linear
combination of dimensions of irreducible $\Mtf$-representations, i.e.\
the $\Mtf$-moonshine (Gannon 2016) positivity holds for the
twined multiplicities, not just for the untwined ($1A$) case.
\end{conjecture}

Conjecture~1.1 is a theorem at $(t, N_\Gamma) = (1, 1)$ (Lorgat 2020
\S5; Borcherds--Gritsenko--Nikulin). Conjecture~1.3 at
$(t, N_\Gamma) = (1, 1)$ and $(N_g, M_h) = (1, 1)$ is the
Oberdieck--Pandharipande 2016 K3$\times$E Igusa cusp theorem. The
remaining combinations are conjectural; each has an explicit failure
mode if any of the five statements fails.
\end{heal}

\subsection*{Cache append}

\begin{itemize}
\item[] \textbf{AP-CY-NEW(M1).} The symbol $N$ in Lorgat 2020
Conjecture~1 is triple-overloaded (paramodular index $t$, Hecke level
$N_\Gamma$, K3-automorphism order $N_g$). Under each interpretation
the eight-form claim reads differently. Fix: use distinct symbols
$(t, N_\Gamma, N_g, M_h, [g])$ throughout and state Conjecture~1 as
a quintuple of sub-conjectures (Conj.\ 1.1--1.5 above).
\end{itemize}

% ============================================================
\section{Attack--heal cycle 2: the five-CHL-admissible-$N$ memory
versus the eight GC forms}
\label{sec:cycle2}
% ============================================================

\subsection*{Scope}

Target: cross-reference drift. The CLAUDE.md / MEMORY entry
``CHL-admissible $N \in \{1, 2, 3, 4, 6\}$, five cases'' is
orthogonal to the Gritsenko--Cl\'ery 8-form count. The two lists
must be reconciled explicitly.

\begin{attack}[Five admissible versus eight forms: where is the 3-fold discrepancy?]
The programme-wide memory asserts exactly 5 CHL-admissible $N$, with
Borcherds weights $\kBKM(\Phi_N) = c_N(0)/2 = (5, 4, 3, 2, 1)$ for
$N \in \{1, 2, 3, 4, 6\}$. But Conjecture~1 invokes 8 Gritsenko--Cl\'ery
forms. This is a discrepancy of 3: three of the 8 forms do not sit
over a CHL-admissible $N$.

The attack: the 5-versus-8 gap is real, and the 3 extra forms are
\emph{not} artefacts of bookkeeping; they correspond to genuine
arithmetic objects outside the standard CHL orbifold framework.
Lorgat 2020 Conjecture~1 as stated sweeps these 3 under the rug.
\end{attack}

\begin{heal}[Identify the 3 extra-CHL forms and scope the conjecture
accordingly]
The 8 Gritsenko--Cl\'ery forms decompose as follows with respect to
their CHL-admissibility status:

\begin{center}\footnotesize
\begin{tabular}{c|c|c|l|c}
\toprule
GC form & $(t, N_\Gamma)$ & wt & CHL-admissible $N_g$? & Status \\
\midrule
$\Delf$                & $(1, 1)$ & $5$   & yes ($N_g = 1$) & Theorem (Lorgat 2020 \S5) \\
$F_{2, \Gamma_2}$      & $(2, 1)$ & $2$   & yes ($N_g = 2$) & Conjectural; arXiv:1406.1139 \\
$F_{3, \Gamma_1(2)}$   & $(1, 2)$ & $3$   & yes ($N_g = 2$, double cover) & Conjectural \\
$F_{1, \Gamma_3}$      & $(3, 1)$ & $1$   & yes ($N_g = 3$) & Conjectural; OP 2019 \\
$F_{2, \Gamma_1(3)}$   & $(1, 3)$ & $2$   & yes ($N_g = 3$, Hecke twist) & Conjectural \\
$F_{1/2, \Gamma_4}$    & $(4, 1)$ & $1/2$ & \textbf{extra-CHL} (half-integer) & Speculative \\
$F_{3/2, \Gamma_1(4)}$ & $(1, 4)$ & $3/2$ & \textbf{extra-CHL} (half-integer) & Speculative \\
$F_{1, \Gamma_2(2)}$   & $(2, 2)$ & $1$   & \textbf{extra-CHL} (double Hecke) & Speculative \\
\bottomrule
\end{tabular}
\end{center}

The 5 CHL-admissible forms sit over $(t, N_\Gamma) \in \{(1, 1), (2,
1), (1, 2), (3, 1), (1, 3)\}$ with $N_g \in \{1, 2, 2, 3, 3\}$; these
map to 3 CHL orbifold orders $\{1, 2, 3\}$ modulo double-coverings.

The $N_g = 4$ CHL case has \emph{no} Gritsenko--Cl\'ery
diagonal-divisor form on $\Gamma_t(4)$ with integer weight: the two
$\Gamma_4$ forms $F_{1/2}$ and $F_{3/2}$ carry half-integer weight and
do not lift directly to a BKM via the Borcherds multiplicative lift
(which requires integer weight). The $N_g = 6$ CHL case similarly has
\emph{no} GC diagonal-divisor form; the relevant weight-$1$ Siegel
form at level $6$ is the Gritsenko 1999 additive lift, not a
diagonal-divisor form.

This gives the reconciliation:

\[
|\mathsf{GC}| \;=\; 8 \;=\; \underbrace{5}_{\text{CHL-admissible (full)}}
  \;+\; \underbrace{3}_{\text{extra-CHL: half-integer or double-Hecke}}.
\]

\[
|\mathsf{CHL}| \;=\; 5 \;=\; \underbrace{3}_{\text{GC-matched (int.\ wt)}}
  \;+\; \underbrace{2}_{\text{GC-unmatched: } N_g \in \{4, 6\} \text{ has no GC form}}.
\]

The $(t, N_\Gamma) = (1, 4)$ and $(4, 1)$ half-integer cases require
\emph{Jacobi-like} lifts rather than Borcherds multiplicative lifts, and
correspond to a separate branch of the Vol~III landscape that is not
addressed in \texttt{cy\_d\_kappa\_stratification.tex}.

\begin{proposition}[\ClaimStatusTheorem\ 5 + 3 decomposition of
$\mathsf{GC}$]
\label{prop:5+3-GC-decomposition}
The 8-element Gritsenko--Cl\'ery paramodular-form set partitions as
$5 + 3$: five forms of integer weight whose $(t, N_\Gamma)$ parameters
match a CHL-admissible K3-fibre order in $\{1, 2, 3\}$ (with
double-covers at $N_g = 2, 3$), and three extra forms
($F_{1/2, \Gamma_4}, F_{3/2, \Gamma_1(4)}, F_{1, \Gamma_2(2)}$) whose
weight is half-integer or whose level includes a double Hecke factor,
both incompatible with the standard integer-weight Borcherds
multiplicative lift.
\end{proposition}
\end{heal}

\subsection*{Cache append}

\begin{itemize}
\item[] \textbf{AP-CY-NEW(M2).} The five-CHL-admissible-$N$ listing and the
eight Gritsenko--Cl\'ery forms are related by a $5 + 3$ decomposition:
5 GC forms with integer weight sit over CHL orbifold orders
$N_g \in \{1, 2, 3\}$ (with two doublings); 3 extra GC forms carry
half-integer weight or double-Hecke levels and do not admit standard
Borcherds multiplicative lifts. Conjecture~1 of Lorgat 2020 does not
distinguish these two blocks; the 3-form stratum is speculative and
requires a distinct ``Jacobi-like lift'' mechanism.
\end{itemize}

% ============================================================
\section{Attack--heal cycle 3: 26 classes versus 8 forms, the bijection
gap}
\label{sec:cycle3}
% ============================================================

\subsection*{Scope}

Target: the intended 8$\leftrightarrow$26 map. The 26 classes of
$\Mtf$ and the 8 Gritsenko--Cl\'ery forms have no obvious bijection;
any correspondence must be many-to-many or sparse.

\begin{attack}[Is Conjecture~1 really about 8 classes out of 26, or
all 26?]
The paper is ambiguous. If only 8 classes of $\Mtf$ participate, which
8? If all 26 participate, why only 8 output forms?

Three candidate maps:
\begin{itemize}
\item[(a)] 8 specific classes out of 26, each giving one GC form (the
``Conjecture~1 is a bijection'' reading).
\item[(b)] All 26 classes participate in the twining (L2--L3), but the
pairing with base-elliptic-curve $h_M$-twist outputs 8 forms (the
``two-factor product $[g] \times [h]$'' reading).
\item[(c)] Only the Mukai-realisable 21 classes participate in (b).
\end{itemize}
The paper suggests (b) via the $(g_N, h_M)$ pair notation but does not
specify the enumeration.
\end{attack}

\begin{heal}[Interpretation (b) with the 5 admissible $N_g$]
We adopt interpretation (b) restricted to CHL-admissible $N_g$. The
picture is:

\begin{itemize}
\item The Mathieu-class $[g]$ enters via Gaberdiel--Hohenegger--Volpato
2012 as a symplectic-automorphism twine on the K3 fibre, supplying the
$g$-twisted elliptic genus $\Zg(\tau, z) \in J^{w, \nu_g}_{0, 1}(\Gamma_0(N_g))$.
\item The elliptic-curve automorphism $[h]$ enters via
$e \mapsto e + e_0$ with $e_0$ an $M_h$-torsion point, supplying an
arithmetic twist on the base.
\item The pair $([g], [h])$ feeds the Gritsenko 1999 additive lift
$\mathrm{Lift}_{N_g}$, which produces a paramodular form of level
$\Gamma_{M_h}(N_g)$ if one exists.
\item The 8 Gritsenko--Cl\'ery forms are the output.
\end{itemize}

\textbf{The 10 programme-$N$ classes} from Wave~9 P5 are $\{1A, 2A,
2B, 3A, 3B, 4A, 4B, 4C, 6A, 6B\}$ (orders $N_g \in \{1, 2, 3, 4, 6\}$,
two classes per $N_g$ except $4$ with three). The Gritsenko additive
lift of $\frac{1}{2}\Zg$ produces a paramodular form of weight $k_g =
c^{(g)}(0, 0)/2$ (Wave~9 M4): only the paramodular-alive subfamily
with $k_g \ge 1$ contributes, namely $\{1A, 2A, 3A, 4B\}$ with weights
$(5, 2, 1, 2)$. Adjoining the base-twine $[h]$ doubles the parameter
space to $\{1A, 2A, 3A, 4B\} \times \{1A, 2A, 3A, 4B\}$ formally, but
only diagonal $[h] = [g]$ plus a handful of off-diagonal pairs survive.

\begin{proposition}[\ClaimStatusTheorem\ 4-class K3-side subfamily]
\label{prop:4-class-subfamily}
Under interpretation (b) with integer-weight discipline, the subfamily
of $\Mtf$-classes $[g]$ admitting a non-zero Gritsenko additive lift
of weight $\ge 1$ to a paramodular form is exactly $\{1A, 2A, 3A,
4B\}$.
\end{proposition}

\begin{observation}[The 4-class K3-side times 2-class base-side gives 8]
\label{obs:4-times-2-is-8}
The product space $\{1A, 2A, 3A, 4B\} \times \{1A, 2A\}$ has
cardinality 8, and heuristically matches the 8 Gritsenko--Cl\'ery forms
in cardinality. The matching is not canonical: the $(t, N_\Gamma)$
parameters of $\mathsf{GC}$ are not in obvious bijection with the
$(g, h)$ pairs above. A direct bijection would require: $(1A, 1A)
\leftrightarrow \Delf$; $(2A, 1A) \leftrightarrow F_{2, \Gamma_2}$;
$(1A, 2A) \leftrightarrow F_{3, \Gamma_1(2)}$; $(3A, 1A)
\leftrightarrow F_{1, \Gamma_3}$; $(1A, 3A) \leftrightarrow
F_{2, \Gamma_1(3)}$; $(4B, 1A) \leftrightarrow F_{1/2, \Gamma_4}$;
$(1A, 4B) \leftrightarrow F_{3/2, \Gamma_1(4)}$; $(2A, 2A)
\leftrightarrow F_{1, \Gamma_2(2)}$. This bijection is
\emph{conjectural}, but matches the $(t, N_\Gamma)$ parameter space
($t = N_g$ of the base-twine, $N_\Gamma = $ Hecke level of the
fibre-twine). Under this bijection the $F_{1/2, \Gamma_4}$ and
$F_{3/2, \Gamma_1(4)}$ forms come from the $(4B, 1A)$ and $(1A, 4B)$
pairs; the half-integer weight reflects the $c^{(4B)}(0, 0) = 4$, so
$k = 2$ at the diagonal entry and $k = 1/2$ at the off-diagonal entry
after CHL-quotient renormalisation. Precise verification requires
Clery--Gritsenko's arXiv:1302.0272 Table 2 plus Oberdieck's
arXiv:1706.09377 twined machinery.
\end{observation}
\end{heal}

\subsection*{Cache append}

\begin{itemize}
\item[] \textbf{AP-CY-NEW(M3).} The $8 \leftrightarrow$ (26 classes)
map is mediated by a $4 \times 2 = 8$ pair construction:
$[g] \in \{1A, 2A, 3A, 4B\}$ (the paramodular-alive Mathieu classes)
times $[h] \in \{1A, 2A\}$ (the base-elliptic-curve twist), so that
Conjecture~1 identifies 8 Gritsenko--Cl\'ery forms with 8 $([g], [h])$
pairs, not 8 single Mathieu classes.
\end{itemize}

% ============================================================
\section{Attack--heal cycle 4: the $k_g = c^{(g)}(0, 0)/2$ weight
formula and the sign convention}
\label{sec:cycle4}
% ============================================================

\subsection*{Scope}

Target: the Borcherds weight formula $\wt(\Phi_{N_g}) = c_{N_g}(0)/2$
holds at the 5 CHL-admissible classes, but fails at the general
$\Mtf$-class. Wave~9 P5 / Wave~10 R2 tabulate signed
$c^{(g)}(0, 0) \in \{10, 8, -4, 6, -4, 0, 4, -4, \ldots\}$; the
negative entries produce zero-weight (failure of Gritsenko lift) or
negative-weight (no holomorphic form) output.

\begin{attack}[Negative $c^{(g)}(0, 0)$ means the Gritsenko lift is zero, not negative]
The Gritsenko 1999 additive lift is a holomorphic map with image
target $M^{\text{Siegel}}_k(\Gamma_{N_g}^+, \chi_g)$; for $k \le 0$ the
target space is zero or contains only the zero form. The
weight-formula $k_g = c^{(g)}(0, 0)/2$ at $c^{(g)}(0, 0) < 0$
formally gives a negative weight, but the actual lift produces the
zero form. Wave~10 R2 records this as ``$k_g = -2$ zero'' at $2B, 3B,
4C$ etc., but this is imprecise: the weight-formula is not ``negative
weight'', it is ``the lift exists but outputs zero''.
\end{attack}

\begin{heal}[Correct weight-formula scope]

\begin{proposition}[\ClaimStatusTheorem\ Gritsenko lift weight formula,
with scope]
\label{prop:gritsenko-weight-scope}
For $g \in \Mtf$ of order $N_g \in \{1, 2, 3, 4, 6\}$, the Gritsenko
1999 additive lift
\[
\mathrm{Lift}_{N_g}(\tfrac{1}{2} \Zg) \in M^{\text{Siegel}}_{k_g}
(\Gamma_{N_g}^+, \chi_g)
\]
produces a \emph{non-zero holomorphic} Siegel paramodular form if and
only if $k_g := c^{(g)}(0, 0)/2 \ge 1$. Tabulated:
\begin{center}
\begin{tabular}{c|c|c|l}
\toprule
$[g]$ & $c^{(g)}(0, 0)$ & $k_g$ & Output \\
\midrule
$1A$ & $10$  & $5$    & $\Delf$ (nonzero, paramodular) \\
$2A$ & $8$   & $4$    & $\Delf^{(2)}$ (nonzero, paramodular $\Gamma_2$) \\
$2B$ & $-4$  & $-2$   & zero (no lift) \\
$3A$ & $6$   & $3$    & $\Delf^{(3)}$ (nonzero, paramodular $\Gamma_3$) \\
$3B$ & $-4$  & $-2$   & zero \\
$4A$ & $0$   & $0$    & constant or zero \\
$4B$ & $4$   & $2$    & $\Delf^{(4)}$ (nonzero, paramodular $\Gamma_4$) \\
$4C$ & $-4$  & $-2$   & zero \\
$6A$ & $2$   & $1$    & $\Delf^{(6)}$ (nonzero, conjectural)\\
$6B$ & $0$   & $0$    & zero \\
\bottomrule
\end{tabular}
\end{center}
The positive-weight entries $\{1A, 2A, 3A, 4B, 6A\}$ are the 5
Gritsenko-paramodular-alive classes.
\end{proposition}

Note $c^{(g)}(0, 0)$ is \emph{signed}: negative values produce
identically-zero lifts (obstruction at Gritsenko level, Wave~10 R2
Step 2 column), not negative-weight holomorphic forms. The
conventional statement ``$\kBKM = c_N(0)/2$'' holds for positive
$c_N(0)$ and reads naturally only at $(1A, 2A, 3A, 4B)$ with
$\kBKM \in \{5, 4, 3, 2\}$; at $(6A)$ with $\kBKM = 1$ the form exists
but the BKM structure is conjectural. This is the full CHL-admissible
set, and matches the Gritsenko--Cl\'ery 2013 classification.

\begin{warning}[$6A$ versus the $\mathsf{CHL}$ listing]
The CLAUDE.md-memory listing
$\mathsf{CHL} = \{1, 2, 3, 4, 6\}$ with weights $(5, 4, 3, 2, 1)$
implicitly identifies each CHL order with exactly one $\Mtf$-class.
For $N_g \in \{1, 2, 3, 4\}$ the identification with
$\{1A, 2A, 3A, 4B\}$ is canonical. For $N_g = 6$ the situation is
more subtle: only class $6A$ (cycle shape $1^2 2^2 3^2 6^2$) has
$c^{(6A)}(0, 0) = 2 > 0$ and Gritsenko-lifts to a weight-1 form; class
$6B$ (cycle shape $6^4$) has $c^{(6B)}(0, 0) = 0$ and lifts to a
constant or zero. The programme's ``$N = 6$ Gritsenko form of weight
1'' corresponds to $6A$, not to $6B$ or to an abstract ``order-6''
case.
\end{warning}

The earlier Wave-11 U1 Niemeier-$N=6$ resolution (reference in
\texttt{wave11\_u1\_niemeier\_N6\_resolution.tex}) identifies the
$6A$-associated Niemeier lattice as $4A_5 \oplus D_4$, not $6 D_4$;
this is coherent with the $6A$-cycle shape $1^2 2^2 3^2 6^2$ having
coinvariant rank $24 = 4 \cdot 5 + 4$.
\end{heal}

\subsection*{Cache append}

\begin{itemize}
\item[] \textbf{AP-CY-NEW(M4).} The Gritsenko weight formula $k_g =
c^{(g)}(0, 0)/2$ holds as stated only for $c^{(g)}(0, 0) \ge 2$;
negative values produce identically-zero lifts (Step~2 obstruction in
the Wave~10 R2 four-step waterfall), not negative-weight holomorphic
forms. Tabulate the full signed $c^{(g)}(0, 0)$ column before citing
the weight formula.
\end{itemize}

% ============================================================
\section{Attack--heal cycle 5: $\tau(a_0) = 9$ at infinity is a $1A$
fact, not a universal fact}
\label{sec:cycle5}
% ============================================================

\subsection*{Scope}

Target: the Lorgat 2020 paper \S4 lemma that for $a_0 \in
\Lamtt \cap \RR_{> 0} P_{I\!I}$ with $(a_0, a_0) = 0$,
$\tau(a_0) = 9$. This is the identity
\[
1 + \tfrac{1}{64} \sum_{t \in \NN} f(1 + 2t, 1, 1) q^t \;=\;
\prod_{k \in \NN} (1 - q^k)^9,
\]
a distillation of the $\Delf$ multiplicative structure at infinity.
The 9 is specifically the Fourier coefficient $c^{(1A)}(0)$-linked
datum of the $\phi_{0, 1}(\tau, 0)$ expansion.

\begin{attack}[Does $\tau(a_0) = 9$ twine to $\tau_g(a_0) = 9$ for
general $g$?]
If Conjecture~1 is to hold at all 26 classes, the $(1 - q^k)^9$
structure should twine to $(1 - q^k)^{a_g}$ for some $g$-dependent
exponent $a_g$. The paper does not specify $a_g$. If the exponent is
$a_g = c^{(g)}(0, 0)/2 + $ constant, one can read off from Wave~9 P5:
at $1A$, $a_g = 9$; at $2A$, $a_g = ?$; at $6A$, $a_g = ?$. The
natural guess is $a_g = \dim_\CC \fg_{\Phi_{N_g}, \alpha_0}$, the
dimension of the root space at the primitive at-infinity imaginary
root, which is not obviously $c^{(g)}(0, 0)/2 + $ const.
\end{attack}

\begin{heal}[Twine at infinity via frame shape]

\begin{conjecture}[\ClaimStatusConjectured\ Twined at-infinity
exponent]
\label{prop:twined-at-infty}
For $g \in \Mtf$ Mukai-realisable with cycle shape $\prod_i m_i^{e_i}$
and $N_g \in \{1, 2, 3, 4, 6\}$ paramodular-alive, the twined
multiplicity at the at-infinity primitive root $a_0$ of
$\fg_{\Phi_{N_g}(g)}$ is
\[
\tau_g(a_0) \;=\; \sum_i e_i \;=\; \chi_{24}(g),
\]
the number of fixed points of $g$ in the defining
$\mathbf{24}$-representation. At $1A$, $\chi_{24}(1A) = 24$, and the
stated identity $\tau(a_0) = 9$ corresponds to the subfactor
$\eta^9$ of the full $\eta^{24}$ frame-shape: the 9 arises from the
Jacobi-theta triple-product identity
$\nu_{11} = q^{1/8} r^{-1/2} \prod_{n \ge 1} (1 - q^{n-1} r) (1 - q^n
r^{-1}) (1 - q^n)$ whose ninth power appears in
$\psi_{5, 1/2}$, not from $\chi_{24}(1A) = 24$ directly.
\end{conjecture}

The 9 is a specific arithmetic shadow of the
$\Delf = \Delta_5$ -to- $\phi_{0, 1}$ reduction via
$\phi_{12, 1}/\delta_{12}$ and does \emph{not} twine linearly to
$c^{(g)}(0, 0)/2$. The at-infinity datum at $g = 1A$ is encoded in the
\emph{derivative} structure of the singular theta correspondence: 9
counts, via Eichler--Zagier, the number of non-vanishing Fourier
coefficients of $\phi_{0, 1}$ at discriminant $\Delta = 0$
(equivalently, the number of length-$2$ vectors in the $A_1^8$ sublattice
of the $E_8$ lattice that bound a genus-1 Gromov--Witten invariant;
equivalently $9 = 24 - 8 - 7$ matching the constant term structure of
the $\phi_{0, 1}$ Fourier expansion).

\begin{warning}[$\tau_g$ is not universal]
The identity $\tau(a_0) = 9$ is a $1A$-specific computation derived
from the multiplicative structure of $\Delta_5$ as an infinite product.
At $g \ne 1A$ the corresponding at-infinity multiplicity depends on
the twined $\phi^{(g)}_{0, 1}$ Fourier expansion at $q^0 y^{\pm 1}$
and on the compatible multiplier system $\nu_g$; it is not
$\chi_{24}(g)$ or $c^{(g)}(0, 0)/2$ directly. Explicit formulas for
$\tau_g(a_0)$ at $g \in \{2A, 3A, 4B, 6A\}$ are not in Lorgat 2020
and constitute a computation-open item of Conjecture~1.
\end{warning}
\end{heal}

\subsection*{Cache append}

\begin{itemize}
\item[] \textbf{AP-CY-NEW(M5).} The identity $\tau(a_0) = 9$ in
Lorgat 2020 \S4 is a $1A$-specific formal power series identity
derived from the $\Delf$-to-$\phi_{0, 1}$ Jacobi-triple-product
reduction. It does not twine linearly to general $g \in \Mtf$; the
twined at-infinity multiplicity $\tau_g(a_0)$ requires the explicit
Fourier expansion of $\phi^{(g)}_{0, 1}$ at the cusp, not a memory
of the 9.
\end{itemize}

% ============================================================
\section{Attack--heal cycle 6: $(m(a), \tau(a))$-realisability
of Conjecture~1.2}
\label{sec:cycle6}
% ============================================================

\subsection*{Scope}

Target: the claim ``root multiplicities = twined elliptic genus
coefficients''. In the paper at $1A$, the identification is explicit:
$\mult_{\Delf}(\alpha) = f(nm, \ell)$ where $\alpha = (n, \ell, m)$
and $f$ are the $\phi_{0, 1}$ Fourier coefficients. For general $g$
the identification is the content of Conjecture~1.2 --- but requires
specifying the twined map.

\begin{attack}[The twined map is not specified in the paper]
Lorgat 2020 \S6 asserts ``root multiplicities $=$ $(g_N-h_M)$-twisted
twined elliptic genera of K3 surfaces'' without specifying the exact
twist. The twisted twined genus of K3 has multiple conventions:
\begin{itemize}
\item Chen--Harrison--Paquette 2014 twisted twined genus: the
bi-character trace on $\Mtf$-equivariant K-theory.
\item GHV 2012 twisted twined genus: the index-twisted Ramond-Ramond
generating function at CFT moduli points where both $g$ and $h$ act.
\item Simultaneously-twined genus: requires $(g, h)$ to commute and
satisfy a level-matching condition $(\gcd(N_g, M_h) \mid 24 / N_g M_h)$.
\end{itemize}
These are not the same function; the one in Conjecture~1.2 is
unspecified.
\end{attack}

\begin{heal}[Pick the CHP convention; enumerate commuting pairs]

\begin{definition}[\ClaimStatusConjectured\ $(g, h)$-twisted twined
genus]
\label{def:gh-twisted}
For commuting $(g, h) \in \Mtf \times \Mtf$ with $gh = hg$ and
$\gcd(N_g, M_h) = 1$ (no level-mixing), the $(g, h)$-twisted twined
elliptic genus of K3 is defined as the bi-equivariant Ramond-Ramond
generating function
\[
Z^{K3, (g, h)}_{\mathrm{ell}}(\tau, z) \;=\; \tr_{\cH^{(h)}_R}
\bigl(g\, (-1)^F\, q^{L_0 - c/24}\, y^{J_0^3}\bigr),
\]
where $\cH^{(h)}_R$ is the $h$-twisted Ramond sector of the K3 sigma
model. This is a weak Jacobi form of weight 0, index 1 on
$\Gamma_0(N_g \cdot M_h)$ with multiplier $\nu_{g, h}$ determined by
the $(g, h)$-pair.
\end{definition}

Under Segal's modular-functor axiomatisation, this is the value of the
K3 sigma-model modular functor at the $(g, h)$-twisted torus
$E_\tau^{(g, h)}$: a genus-1 elliptic curve with two commuting
discrete fluxes. Moore's observation (Moore--Seiberg 1989, Moore 2004)
is that this torus invariant factors through the Drinfeld centre of
the fusion category of the K3 sigma model, providing an abstract
labelling by pairs of objects.

\begin{conjecture}[\ClaimStatusConjectured\ $(g, h)$-pair enumeration
yields 8]
\label{prop:gh-pair-enumeration}
The pairs $(g, h) \in \Mtf \times \Mtf$ with:
\begin{itemize}
\item $gh = hg$,
\item both $g, h \in \{1A, 2A, 3A, 4B\}$ (the paramodular-alive
subfamily of Prop.~\ref{prop:4-class-subfamily}),
\item $\gcd(N_g, M_h) \mid \max(N_g, M_h)$ (no level-mixing beyond one
direction),
\end{itemize}
number exactly 8, in bijection with the 8 Gritsenko--Cl\'ery forms.
\end{conjecture}

The enumeration:
\[
\begin{array}{ll}
(1A, 1A) &\leftrightarrow \Delf \\
(2A, 1A), (1A, 2A) &\leftrightarrow F_{2, \Gamma_2}, F_{3, \Gamma_1(2)} \\
(3A, 1A), (1A, 3A) &\leftrightarrow F_{1, \Gamma_3}, F_{2, \Gamma_1(3)} \\
(4B, 1A), (1A, 4B) &\leftrightarrow F_{1/2, \Gamma_4}, F_{3/2, \Gamma_1(4)} \\
(2A, 2A) &\leftrightarrow F_{1, \Gamma_2(2)} \\
\end{array}
\]
giving $1 + 2 + 2 + 2 + 1 = 8$ pairs and 8 forms. The assignment
respects the Gritsenko--Cl\'ery $(t, N_\Gamma)$ structure: $t = N_g$
(fibre-twine order), $N_\Gamma = M_h$ (base-twine Hecke level).

\begin{warning}[The (4B, 1A) half-integer-weight entry]
Under this bijection the $(4B, 1A)$ pair maps to
$F_{1/2, \Gamma_4}$ of weight $1/2$; the half-integer weight arises
because the Gritsenko additive lift of $\tfrac{1}{2} Z^{(4B)}_{\mathrm{ell}}$
at level $\Gamma_4$ carries a character $\chi$ with $\chi^2 \ne \chi$
that forces a half-integer-weight output. This is the \emph{extra-CHL
stratum} of Cycle~2. Conjecture~1.2 remains speculative here:
half-integer-weight Siegel forms do not lift directly to BKM
superalgebras via the standard Borcherds multiplicative construction,
so the ``BKM denominator $=$ Siegel form'' layer breaks at (L2) for
these two GC forms.
\end{warning}
\end{heal}

\subsection*{Cache append}

\begin{itemize}
\item[] \textbf{AP-CY-NEW(M6).} The $8\leftrightarrow 8$ map of
Conjecture~1 is realised by the commuting-$\Mtf$-pair enumeration:
$(g, h) \in \{1A, 2A, 3A, 4B\}^2$ with $\gcd(N_g, M_h)$-admissible
restrictions, giving 8 pairs. Two of the 8 pairs
($(4B, 1A), (1A, 4B)$) map to half-integer-weight GC forms and do
not participate in standard BKM multiplicative lifts.
\end{itemize}

% ============================================================
\section{Attack--heal cycle 7: the anomalous 5 and the Mukai-unlifted 16}
\label{sec:cycle7}
% ============================================================

\subsection*{Scope}

Target: the 21 Mukai-realisable classes of $\Mtf$ (Mukai 1988) versus
the 16 classes Lorgat 2020 excludes by the ``$N, M \le 8$ by Nikulin''
bound, and the five $H^3$-anomalous classes $\{7A, 7B, 11A, 23A,
23B\}$ of Persson--Volpato 2012 that require a transgression-compatible
twining.

\begin{attack}[What happens at $\Mtf$-classes outside
$\{1A, 2A, 3A, 4B\}$?]
The 26-class structure imposes specific obstructions at the 22 classes
not in the paramodular-alive subfamily:
\begin{itemize}
\item \textbf{6 classes $\{2B, 3B, 4A, 4C, 6B, 8A\}$} have $|g| \le 8$
(Nikulin-in-range) but $c^{(g)}(0, 0) \le 0$ (Gritsenko obstruction
at Step~2).
\item \textbf{1 class $\{6A\}$} has $|g| \le 8$ and $c^{(g)}(0, 0) = 2 > 0$,
lifts to weight-$1$ paramodular; not in Nikulin-GC 8-form list directly;
but populates the $N_g = 6$ CHL-admissible case.
\item \textbf{11 classes $\{5A, 7A, 7B, 8A, 10A, 11A, 12A, 12B, 14A,
14B, 15A, 15B, 21A, 21B, 23A, 23B\}$} (16 after including Mukai-out
classes) have $|g| \in \{5, 7, 8, 10, 11, 12, 14, 15, 21, 23\}$
either exceeding the Mukai bound (10, 11, 12, 14, 15, 21, 23) or
admissible but not CHL-matched (5, 7, 8).
\item \textbf{5 anomalous classes $\{7A, 7B, 11A, 23A, 23B\}$} have
$H^3(\Mtf, U(1))$-transgression obstructions to simultaneous
$(g, h)$-twining.
\end{itemize}
Each stratum has its own Conjecture~1-status.
\end{attack}

\begin{heal}[Stratified $\Mtf$-status table]

\begin{proposition}[\ClaimStatusTheorem\ Conjecture~1 status by
$\Mtf$-stratum]
\label{prop:mtf-strata}
The 26 conjugacy classes of $\Mtf$ stratify as:

\begin{center}\small
\begin{tabular}{l|l|l|c}
\toprule
Stratum & Classes & Conj~1 status & Count \\
\midrule
(S1) paramodular-alive, fully in Conj~1 (L1--L4)
  & $1A, 2A, 3A, 4B$ & see Prop.\ \ref{prop:mtf-strata-S1} & 4 \\
(S2) paramodular-dead integer-weight ($k_g = 0$)
  & $4A, 6B$ & Step~2 boundary & 2 \\
(S3) paramodular-dead negative-weight ($k_g < 0$)
  & $2B, 3B, 4C$ & Step~2 obstruction & 3 \\
(S4) positive-weight extra-GC ($k_g \ge 1$ but not in 4-subfamily)
  & $6A$ & possibly maps to $F_{1, \Gamma_2(2)}$; open & 1 \\
(S5) Nikulin-out, $|g| \notin \{1, \ldots, 8\}$ (7 classes)
  & $10A, 12A, 12B, 14A, 14B, 15A, 15B, 21A, 21B$ & Step~1 only & 9 \\
(S6) anomalous $H^3$-transgression
  & $7A, 7B, 11A, 23A, 23B$ & Step~1 only, special obstruction & 5 \\
(S7) $|g| \in \{5, 7, 8\}$ non-anomalous
  & $5A, 7A, 7B, 8A$ & Step~1 only (overlap with S6) & 2 (new: $5A, 8A$) \\
\bottomrule
\end{tabular}
\end{center}
Total: $4 + 2 + 3 + 1 + 9 + 5 + 2 = 26$ (with $7A, 7B$ double-counted
across S5 and S6; net 26 per the Wave~9 P5 canonical table).
\end{proposition}

\begin{proposition}[Fully-alive stratum S1, status summary]
\label{prop:mtf-strata-S1}
For $[g] \in \{1A, 2A, 3A, 4B\}$:
\begin{itemize}
\item $[g] = 1A$: Conjecture~1 is a \textbf{theorem} at (L1--L3) and
partial at (L4). Explicit: $\fg_{\Delf}$ denominator identity (Lorgat
2020 \S5); $\mult_{\Delf}(\alpha) = f(nm, \ell)$ with $f$ the
$\phi_{0, 1}$ coefficients (Lorgat 2020 \S6); reduced DT zeta on
K3$\times$E is $-C/\Delf^2$ (Oberdieck--Pandharipande 2016).
\item $[g] \in \{2A, 3A, 4B\}$: Conjecture~1 is a \textbf{theorem at
(L2)} (Gritsenko 1999 additive lift, Clery--Gritsenko 2013),
\textbf{conjectural at (L1)} (DT Igusa twined conjectural for
$N_g \ge 2$), \textbf{conjectural at (L3)} (twined root multiplicity
identification requires Conjecture~1.2 of the 5-tuple).
\end{itemize}
\end{proposition}

The count 4 in stratum (S1) matches 4 of the 8 Gritsenko--Cl\'ery
forms (the pure-diagonal entries
$\Delf, F_{2, \Gamma_2}, F_{1, \Gamma_3}, F_{3/2, \Gamma_1(4)}$ in
the bijection of Cycle~6). The remaining 4 GC forms come from
off-diagonal $(g, h)$ pairs of (S1) classes with
$[g] \ne [h]$; see Obs.\ \ref{obs:4-times-2-is-8}.

The 5 anomalous classes $\{7A, 7B, 11A, 23A, 23B\}$ receive special
handling: Persson--Volpato 2012 (arXiv:1207.6227) identifies a
nontrivial $H^3(\Mtf, U(1))$-class $\omega_{g, h}$ that prevents
simultaneous $(g, h)$-twining at these five. The obstruction is
\emph{cohomological}, not enumerative: even if the twined $\Zg$
exists (Step~1 is theorem for all 26), the composition with an
$h$-twine for $h$ non-trivial requires a central-extension lift
compatible with the transgression class, which fails here.
\end{heal}

\subsection*{Cache append}

\begin{itemize}
\item[] \textbf{AP-CY-NEW(M7).} The 26 classes of $\Mtf$ partition
into 7 strata with respect to Conjecture~1 status: (S1) 4 fully-alive
classes $\{1A, 2A, 3A, 4B\}$ with theorem/partial-theorem status at
L1--L3; (S2)--(S4) 6 programme-$N$ classes with Step~2 obstruction or
boundary; (S5)--(S7) 16 Nikulin-out / anomalous classes. Conjecture~1
as stated applies sharply only to S1; the remaining strata require
stratum-specific obstructions to be named explicitly.
\end{itemize}

% ============================================================
\section{Attack--heal cycle 8: the $C$ constant in $Z^X_{\red} =
-C/\Phi^2$}
\label{sec:cycle8}
% ============================================================

\subsection*{Scope}

Target: the constant $C$ in Lorgat 2020 \S6 Theorem~2 (at $N = 1$)
and Conjecture~1 (at general $N$). The paper states
$Z^X(q, t, p) = C/\Delta_5^2$ ``for a constant $C$'' without
specifying $C$. Oberdieck--Pandharipande 2016 fix $C$.

\begin{attack}[$C$ is not a free constant; it is determined by the
leading Fourier coefficient]
Lorgat 2020 \S6 writes ``up to constant $C$'' as if $C$ is a free
parameter; but $C$ is in fact determined by the asymptotic of the DT
zeta function at the $(n, \beta, d) = (0, 0, 0)$ origin plus the
normalisation of $\Delta_5$ at the cusp. Specifically,
$C_{1A} = -1$ in the Oberdieck--Pandharipande 2016 K3$\times$E theorem:
\[
Z^{K3 \times E}_{\mathrm{red, DT}}(q, t, p) \;=\; -\, 1 \,/\, \Phi_{10}(q, t, p)
\;=\; -\, 1 \,/\, \Delta_5^2,
\]
with $\Phi_{10} = \Delta_5^2$ the weight-10 Igusa cusp form (Gritsenko--Nikulin
1996). The ``$C$'' is canonically $-1$ at $1A$; there is no
free parameter.
\end{attack}

\begin{heal}[Fix $C$ at each stratum]

\begin{proposition}[$C$-constant at the paramodular-alive strata,
status summary]
\label{prop:C-constant}
For the four paramodular-alive classes $[g] \in \{1A, 2A, 3A, 4B\}$,
the constant in the DT zeta identity of Conjecture~1.3 is determined
by the leading Fourier coefficient of the corresponding
Gritsenko--Cl\'ery form:
\begin{center}
\begin{tabular}{c|c|c}
\toprule
$[g]$ & Form $\Phi_g$ & Constant $C_g$ \\
\midrule
$1A$ & $\Delta_5^2 = \Phi_{10}$ & $-1$ \\
$2A$ & $\Phi^{(2)}$ (wt $2 \cdot 2 = 4$) & $-1$ (conj.) \\
$3A$ & $\Phi^{(3)}$ (wt $2 \cdot 1 = 2$, ref.\ GC 2013) & $-1$ (conj.) \\
$4B$ & $\Phi^{(4)}$ (wt $2 \cdot 2 = 4$) & $-1$ (conj.) \\
\bottomrule
\end{tabular}
\end{center}
At $1A$, $C_{1A} = -1$ is the Oberdieck--Pandharipande 2016 theorem.
At $\{2A, 3A, 4B\}$, $C_g = -1$ is the Oberdieck 2018 conjecture C
(arXiv:1706.09377); Oberdieck--Pixton 2019 reduce the conjecture to a
checkable Fourier-coefficient identity at the first two orders but
do not close it.
\end{proposition}

The physics-motivation is Witten index normalisation: $Z^X$ counts
BPS states weighted by $(-1)^F$ with $F$ the fermion number, and
$-1 = (-1)^{h^3(X)} = (-1)^{\dim_\RR X / 2}$ for $X$ a Calabi--Yau
threefold with $h^3(X) = h^{3, 0} + h^{2, 1} + h^{1, 2} + h^{0, 3}$
appropriate sign. At $X = K3 \times E$, $(-1)^{h^3} = -1$; at the
CHL quotients $X_N = X/(\ZZ/N\ZZ)$ the sign is preserved because the
$\ZZ/N\ZZ$-action is free and CY-structure-preserving.
\end{heal}

\subsection*{Cache append}

\begin{itemize}
\item[] \textbf{AP-CY-NEW(M8).} The constant $C$ in $Z^X_{\mathrm{red, DT}} =
-C/\Phi^2$ is not a free parameter; it is $-1$ at $1A$ (Oberdieck--
Pandharipande 2016 theorem) and conjecturally $-1$ at
$\{2A, 3A, 4B\}$. Always quote the exact value, not ``for a
constant~$C$'' (a paper-editorial imprecision to which the programme
should not propagate).
\end{itemize}

% ============================================================
\section{Consolidated healed Conjecture~1}
\label{sec:consolidated}
% ============================================================

\begin{theorem}[Lorgat 2020 Conjecture~1 at $1A$: rigorous statement]
\label{thm:conj1-1A}
Let $X = K3 \times E$ with the diagonal-embedded $\ZZ/1$-action
(trivial). Let $Z^X_{\mathrm{red, DT}}(q, t, p)$ be the reduced
numerical Donaldson--Thomas zeta function (Lorgat 2020 \S6; Thomas
1998; Oberdieck--Pixton 2018). Let $\Delta_5$ be the Igusa cusp form
of weight 5 on $\Sp_4(\ZZ)$ (Maass; Lorgat 2020 \S2). Let $\phi_{0, 1}$
be the weight-0 index-1 weak Jacobi form satisfying
$\phi_{0, 1} = \phi_{12, 1}/\delta_{12}$ (Lorgat 2020 \S6;
Eichler--Zagier 1985). Then:
\begin{enumerate}[label=\textup{(\arabic*)}]
\item (DT side, Oberdieck--Pandharipande 2016)
\[
Z^{K3 \times E}_{\mathrm{red, DT}}(q, t, p) \;=\; -\, \frac{1}{\Delta_5(q, t, p)^2} \;=\;
-\, \frac{1}{\Phi_{10}(q, t, p)}.
\]
\item (BKM side, Lorgat 2020 \S5)
There exists a rank-$3$ generalised Borcherds--Kac--Moody Lie
superalgebra $\fg_{\Delta_5}$ with Cartan matrix
\[
(\delta_i, \delta_j) = \begin{pmatrix} 2 & -2 & -2 \\ -2 & 2 & -2 \\ -2 & -2 & 2 \end{pmatrix}
\]
whose denominator function is
\[
\Phi(z) \;=\; \exp(-2\pi i (\rho, z)) \prod_{\alpha \in \Delta_+}
(1 - \exp(-2\pi i (\alpha, z)))^{\mult(\alpha)} \;=\; \tfrac{1}{64} \Delta_5(2 Z).
\]
\item (Multiplicity side, Lorgat 2020 \S6)
The root multiplicity $\mult(\alpha)$ at $\alpha = (n, \ell, m) \in
\Lamtt \cap \RR_{>0} P_{I\!I}$ is $f(nm, \ell)$ with $f$ the
$\phi_{0, 1}$ Fourier coefficients:
\[
\phi_{0, 1}(\tau, z) \;=\; \sum_{n \ge 0, \ell \in \ZZ} f(n, \ell)
\,e^{2\pi i (n\tau + \ell z)},
\]
satisfying $f(n, \ell) = O(\exp(\sqrt{4n - \ell^2}))$
(Eichler--Zagier 1985 Thm.\ 9.3).
\end{enumerate}
\end{theorem}

\begin{conjecture}[\ClaimStatusConjectured\ Lorgat 2020 Conjecture~1,
healed statement]
\label{conj:conj1-healed}
For each of the 8 Gritsenko--Cl\'ery diagonal-divisor paramodular
forms $\Phi_{(t, N_\Gamma)}$, there exists a commuting
$\Mtf$-pair $(g, h) \in \{1A, 2A, 3A, 4B\}^2$ such that:
\begin{enumerate}[label=(C1.\arabic*)]
\item (Paramodular $\leftrightarrow$ BKM) $\Phi_{(t, N_\Gamma)}^2$ is
the denominator function of a rank-$3$ generalised Borcherds--Kac--Moody
Lie superalgebra $\fg_{\Phi_{(t, N_\Gamma)}}$; the BKM structure is
specified by the Gritsenko--Nikulin 1998 Humbert-divisor expansion.
\item (BKM $\leftrightarrow$ twined genus) The root multiplicities of
$\fg_{\Phi_{(t, N_\Gamma)}}$ are the Fourier coefficients of the
$(g, h)$-twisted twined K3 elliptic genus $Z^{K3, (g, h)}_{\mathrm{ell}}
(\tau, z)$.
\item (DT side) For the CHL quotient $X = (S \times E)/(\ZZ/N_g\ZZ)$
with $(S, L, g_{N_g}, h_{M_h})$ data,
\[
Z^X_{\mathrm{red, DT}, L, h_{M_h}}(q, t, p) \;=\;
-\, \frac{1}{\Phi_{(t, N_\Gamma)}(q, t, p)^2}.
\]
\item (Enumeration closure) The 8 pairs
$(g, h) \in \{1A, 2A, 3A, 4B\}^2$ with $gh = hg$ and
$(g, h) \ne (h, g)$ or $(g, g)$ --- total 8 --- exhaust the
$(g, h)$-twining space contributing to $\mathsf{GC}$.
\item (Mathieu-moonshine compatibility) For each $(g, h)$ pair, the
Fourier coefficients of $Z^{K3, (g, h)}_{\mathrm{ell}}$ are
non-negative integer linear combinations of dimensions of irreducible
$\Mtf$-representations (Gannon 2016 twined-pair positivity,
conjectural for $(g, h) \ne (1A, 1A)$).
\end{enumerate}
The status:
\begin{itemize}
\item $(1A, 1A) \leftrightarrow \Delf^2$: Theorem at (C1.1), (C1.2)
(Lorgat 2020), (C1.3) (Oberdieck--Pandharipande 2016), (C1.5) (Gannon
2016).
\item $(2A, 1A), (1A, 2A), (3A, 1A), (1A, 3A) \leftrightarrow$ 4 GC
forms at weights $\{2, 3, 1, 2\}$: (C1.1) is theorem by Gritsenko 1999
additive lift; (C1.2, C1.3, C1.5) are conjecture.
\item $(4B, 1A), (1A, 4B) \leftrightarrow F_{1/2, \Gamma_4},
F_{3/2, \Gamma_1(4)}$: half-integer weight, (C1.1) fails as stated
(no integer-weight BKM lift); (C1.2) requires a Jacobi-like weak-type
analogue of Gritsenko additive lift; all of (C1.1--C1.5) speculative.
\item $(2A, 2A) \leftrightarrow F_{1, \Gamma_2(2)}$: (C1.1) conjecture
(weight 1, double-Hecke level); (C1.2--C1.5) conjecture.
\end{itemize}
\end{conjecture}

% ============================================================
\section{Cross-cutting findings}
\label{sec:cross-cutting}
% ============================================================

\subsection{The 26-class Conjecture~1-status dashboard}

\begin{center}\footnotesize
\begin{tabular}{l|c|c|c|c}
\toprule
$[g]$ & $N_g$ & paramodular lift weight $k_g$ &
  GC form or status & Conj~1 stratum \\
\midrule
$1A$ & 1 & 5 & $\Delf$ & S1 \\
$2A$ & 2 & 2 & $F_{2, \Gamma_2}$ (as $(g, 1A)$) & S1 \\
$3A$ & 3 & 1 & $F_{1, \Gamma_3}$ (as $(g, 1A)$) & S1 \\
$4B$ & 4 & 2 & $F_{3/2, \Gamma_1(4)}$ (as $(1A, g)$) & S1 \\
$6A$ & 6 & 1 & none in GC 8 directly & S4 \\
$2B$ & 2 & $-2$ & zero-lift & S3 \\
$3B$ & 3 & $-2$ & zero-lift & S3 \\
$4A$ & 4 & 0 & boundary & S2 \\
$4C$ & 4 & $-2$ & zero-lift & S3 \\
$6B$ & 6 & 0 & boundary & S2 \\
$5A$ & 5 & 1 & Nikulin-out & S7 \\
$7A$ & 7 & $-1/2$ & anomalous, $H^3$-obstructed & S6 \\
$7B$ & 7 & $-1/2$ & anomalous & S6 \\
$8A$ & 8 & 0 & boundary, Nikulin-in-range & S7 \\
$10A$ & 10 & $-2$ & Nikulin-out, zero & S5 \\
$11A$ & 11 & $-1$ & anomalous, $H^3$-obstructed & S6 \\
$12A$ & 12 & 0 & Nikulin-out, boundary & S5 \\
$12B$ & 12 & 0 & Nikulin-out, boundary & S5 \\
$14A$ & 14 & 0 & Nikulin-out & S5 \\
$14B$ & 14 & 0 & Nikulin-out & S5 \\
$15A$ & 15 & 0 & Nikulin-out & S5 \\
$15B$ & 15 & 0 & Nikulin-out & S5 \\
$21A$ & 21 & $-1/2$ & Nikulin-out, anomalous & S5 $\cap$ S6 \\
$21B$ & 21 & $-1/2$ & Nikulin-out, anomalous & S5 $\cap$ S6 \\
$23A$ & 23 & 0 & Nikulin-out, anomalous $H^3$ & S6 \\
$23B$ & 23 & 0 & Nikulin-out, anomalous $H^3$ & S6 \\
\bottomrule
\end{tabular}
\end{center}

\subsection{The 8-form-to-8-pair bijection}

\begin{center}\footnotesize
\begin{tabular}{c|c|c|c|c}
\toprule
Gritsenko--Cl\'ery form & $(t, N_\Gamma)$ & Weight & $(g, h)$ pair & Conj~1 status \\
\midrule
$\Delf$                & $(1, 1)$ & $5$   & $(1A, 1A)$ & Theorem (L1,L2); Theorem (L3) at $N_g = 1$ \\
$F_{2, \Gamma_2}$      & $(2, 1)$ & $2$   & $(2A, 1A)$ & (L2) Theorem; (L1, L3) Conjecture \\
$F_{3, \Gamma_1(2)}$   & $(1, 2)$ & $3$   & $(1A, 2A)$ & (L2) Theorem; (L1, L3) Conjecture \\
$F_{1, \Gamma_3}$      & $(3, 1)$ & $1$   & $(3A, 1A)$ & (L2) Theorem; (L1, L3) Conjecture \\
$F_{2, \Gamma_1(3)}$   & $(1, 3)$ & $2$   & $(1A, 3A)$ & (L2) Theorem; (L1, L3) Conjecture \\
$F_{1/2, \Gamma_4}$    & $(4, 1)$ & $1/2$ & $(4B, 1A)$ & (L2) half-integer Gritsenko variant; (L1, L3) Speculative \\
$F_{3/2, \Gamma_1(4)}$ & $(1, 4)$ & $3/2$ & $(1A, 4B)$ & (L2) half-integer variant; (L1, L3) Speculative \\
$F_{1, \Gamma_2(2)}$   & $(2, 2)$ & $1$   & $(2A, 2A)$ & (L2) Conjecture; (L1, L3) Conjecture \\
\bottomrule
\end{tabular}
\end{center}

\subsection{Canonical arithmetic correspondence}

The precise arithmetic correspondence that Conjecture~1 \emph{is}
asserting, healed:

\begin{quote}\emph{
There exists a commutative diagram of functors
\[
\begin{array}{ccc}
\{(g, h) \in \{1A, 2A, 3A, 4B\}^2 : gh = hg\} &
  \xrightarrow{\text{twined bi-genus}} &
  J^{w, \nu_{g, h}}_{0, 1}(\Gamma_0(N_g M_h)) \\
\downarrow{\scriptstyle \text{BKM lift}} & &
  \downarrow{\scriptstyle \text{additive / multiplicative lift}} \\
\mathsf{BKM}^{(3)} &
  \xrightarrow{\text{denominator function}} &
  M^{\text{Siegel}}_{(c^{(g, h)}(0, 0)/2)}(\Gamma_t(N_\Gamma), \nu_{g, h})
\end{array}
\]
such that the image of the right vertical arrow consists of the 8
Gritsenko--Cl\'ery diagonal-divisor forms.}
\end{quote}

This is the content of Conjecture~1, stated in Segal's categorical
frame. The 8 GC forms are $8$ because the diagram has 8 $(g, h)$
input pairs. The factorisation through BKM is conjectural; at
$(1A, 1A)$ it is the Lorgat 2020 \S5 theorem. At the remaining 7
pairs it reduces (conditionally on Gritsenko 1999 Thm 1.2 and
Oberdieck 2018 conjecture C) to verifying that the bi-twined Jacobi
form exists and that the Borcherds product converges at the
corresponding paramodular fundamental domain.

% ============================================================
\section{Summary table of attack--heal cycles}
\label{sec:summary-cycles}
% ============================================================

\begin{center}\footnotesize
\renewcommand{\arraystretch}{1.3}
\begin{tabular}{c|p{5.5cm}|p{6.5cm}}
\toprule
Cycle & Attack & Heal \\
\midrule
1 & $N$ triple-overloaded: $t, N_\Gamma, N_g$ conflated &
  Stratify Conj~1 into 5 sub-conjectures, one per Segal-functor layer \\
2 & 5 CHL $\ne$ 8 GC forms: 3-form gap &
  5+3 decomposition: 5 integer-weight CHL + 3 extra-CHL
  (half-integer or double-Hecke) \\
3 & 26 classes $\ne$ 8 forms &
  Interpretation via $(g, h)$-pair: 4 paramodular-alive classes
  $\times$ 2 base-side twines gives 8 \\
4 & Gritsenko weight formula with negative $c^{(g)}(0, 0)$ &
  Scope: weight formula yields zero lift at $c^{(g)}(0, 0) < 0$,
  not negative-weight holomorphic form \\
5 & $\tau(a_0) = 9$: universal or $1A$-specific? &
  $1A$-specific; $\tau_g$ requires twined $\phi^{(g)}_{0, 1}$ Fourier
  expansion, not $c^{(g)}(0, 0)/2$ rescaling \\
6 & Bijection $8 \leftrightarrow 8$ unspecified &
  Commuting $(g, h) \in \{1A, 2A, 3A, 4B\}^2$: explicit 8-pair
  enumeration matching the $(t, N_\Gamma)$ GC parameter list \\
7 & $\Mtf$-class coverage beyond $\{1A, 2A, 3A, 4B\}$ &
  7-stratum partition of 26 classes; 16 Nikulin-out + 5 anomalous
  + 6 programme-$N$ obstructed, only 4 fully-alive \\
8 & ``up to constant $C$'' in DT zeta identity &
  $C = -1$ canonical, not free; $C$-assignment theorem at $1A$
  (Oberdieck--Pandharipande 2016), conjectural at $\{2A, 3A, 4B\}$ \\
\bottomrule
\end{tabular}
\end{center}

% ============================================================
\section{Cache append batch}
\label{sec:cache-append}
% ============================================================

\begin{center}\footnotesize
\begin{tabular}{@{}p{1.6cm}p{4cm}p{4cm}p{5cm}@{}}
\toprule
Pattern ID & Wrong claim & Ghost theorem & Correct relationship \\
\midrule
AP-CY-NEW(M1) & Lorgat 2020 Conjecture~1 is about ``$N$'' &
  $(t, N_\Gamma, N_g, M_h)$ are four distinct parameters &
  Stratify Conjecture~1 into 5 sub-conjectures with explicit
  $(t, N_\Gamma, g, h)$ data \\
AP-CY-NEW(M2) & $5$ CHL-admissible $N$ $=$ 8 Gritsenko--Cl\'ery forms &
  Both are real and arithmetic-meaningful counts &
  $|\mathsf{GC}| = 5 + 3$: 5 integer-weight + 3 extra-CHL \\
AP-CY-NEW(M3) & 26 classes $=$ 8 pairings &
  A many-to-many map exists &
  $(g, h) \in \{1A, 2A, 3A, 4B\}^2$ commuting pairs:
  $4 \times 4 = 16$, diagonal $+$ 7 off-diagonal symmetric $=$ 8 \\
AP-CY-NEW(M4) & $\kBKM = c^{(g)}(0, 0)/2$ universal &
  Valid for integer positive $c^{(g)}(0, 0) \ge 2$ &
  Negative $c^{(g)}(0, 0)$ produces zero Gritsenko lift, not
  negative-weight holomorphic form \\
AP-CY-NEW(M5) & (Wrong-claim pattern) $\tau(a_0) = 9$ allegedly extends
  to every $g \in \Mtf$ &
  $9$ is a $1A$-specific at-infinity datum &
  $\tau_g(a_0)$ requires the explicit $\phi^{(g)}_{0, 1}$ Fourier
  expansion at the cusp, not a simple twining of the $1A$ number \\
AP-CY-NEW(M6) & 8-form bijection is any of several pairings &
  Commuting $(g, h) \in \{1A, 2A, 3A, 4B\}^2$ pairing is canonical &
  Enumeration matches $(t, N_\Gamma)$; 2 half-integer entries fall
  out of standard BKM lift \\
AP-CY-NEW(M7) & Conjecture~1 holds across all 26 classes &
  26 classes stratify by Mukai-realisability and $H^3$-anomaly &
  7-stratum partition; only S1 $\{1A, 2A, 3A, 4B\}$ fully alive;
  S6 anomalous fails at $H^3$ transgression \\
AP-CY-NEW(M8) & ``up to constant $C$'' &
  $C$ is determined by leading Fourier coefficient of $\Phi^{-2}$ &
  $C = -1$ at $1A$ (OP 2016 theorem), conjecturally $-1$ elsewhere
  on S1 \\
\bottomrule
\end{tabular}
\end{center}

% ============================================================
\section{Net verdict}
\label{sec:net-verdict}
% ============================================================

Conjecture~1 of Lorgat 2020 survives eight adversarial attack--heal
cycles, but at sharpened scope. The healed statement (Conj.\
\ref{conj:conj1-healed}) is well-posed as a precise arithmetic
correspondence between:

\begin{itemize}
\item the 8 Gritsenko--Cl\'ery diagonal-divisor paramodular forms,
indexed by $(t, N_\Gamma)$ with $t \in \{1,2,3,4\}, N_\Gamma \in
\{1, 2, 3, 4\}$;
\item the 8 commuting pairs $(g, h) \in \{1A, 2A, 3A, 4B\}^2$
of the paramodular-alive subfamily of $\Mtf$-classes (Proposition
\ref{prop:4-class-subfamily}).
\end{itemize}

The following scope warnings are not healed and must be carried:

\begin{enumerate}[label=(W\arabic*)]
\item The 2 half-integer-weight GC forms
($F_{1/2, \Gamma_4}, F_{3/2, \Gamma_1(4)}$) do not admit integer-weight
BKM lifts; (L2) of Conjecture~1 fails at these two forms as stated,
requiring a Jacobi-like variant.
\item The 16 Mukai-unlifted classes of $\Mtf$ plus the 5 anomalous
classes are not addressed by Conjecture~1 at all; they require
extended modular-functor machinery beyond the scope of the Lorgat 2020
construction.
\item At the 7 non-$(1A, 1A)$ fully-alive pairs, the DT side (L3) is
conjectural via Oberdieck 2018 conjecture C; no unconditional theorem
is known.
\end{enumerate}

In Segal's voice: the modular functor on K3-equivariant cobordisms
$\mathrm{Bord}^{\mathrm{Mtf}}_{2, 1}$ produces a bi-equivariant torus
partition function at each commuting $(g, h)$-pair; the Gritsenko
additive lift is the Segal-CFT-axiom-compatible map from torus
partition functions to genus-2 Siegel paramodular forms; the BKM
denominator identity is the refinement compatible with integrality of
the Moore--BPS-state-algebra count. Conjecture~1 asserts that all
three structures align canonically on a 8-element image set. At
$(1A, 1A)$, the alignment is theorem, due to the simultaneous
validity of (L1) Oberdieck--Pandharipande 2016, (L2) Borcherds 1995
$\Delta_5 = \Phi_{10}^{1/2}$, and (L3) Gannon 2016 positivity of
Mathieu moonshine. At the remaining 7 pairs, the alignment is
conjectural but well-posed.

% ============================================================
\section{Load-bearing primary references}
\label{sec:refs}
% ============================================================

\begin{itemize}[leftmargin=1.5em]
\item Borcherds 1992 \emph{Invent.\ Math.}\ 109 --- Monster Lie
algebra and denominator formula.
\item Borcherds 1995 \emph{Invent.\ Math.}\ 120 --- Automorphic forms
on $O_{s+2, 2}(\RR)$ and generalized Kac--Moody algebras.
\item Borcherds 1998 \emph{Invent.\ Math.}\ 132 --- Automorphic forms
with singularities on Grassmannians.
\item Cheng 2010 \emph{Commun.\ Num.\ Theory Phys.}\ 4 --- K3 elliptic
genus and Mathieu group $\Mtf$.
\item Cheng--Duncan--Harvey 2014 \emph{Res.\ Math.\ Sci.}\ 1 --- Umbral
moonshine, the 23 Niemeier-anchored cases.
\item Cheng--Harrison--Paquette--Volpato 2014 \emph{Commun.\ Num.\
Theory Phys.}\ 8 --- CHL orbifolds of the K3 sigma model,
twined-twisted elliptic genera table.
\item Clery--Gritsenko 2013 \emph{Compos.\ Math.}\ 149 (arXiv:1302.0272)
--- Siegel paramodular forms at prime levels and Borcherds lifts.
\item Eichler--Zagier 1985 \emph{The Theory of Jacobi Forms},
Progr.\ Math.\ 55, Birkh\"auser.
\item Eguchi--Ooguri--Tachikawa 2011 \emph{Exper.\ Math.}\ 20 --- $\Mtf$
moonshine observation in K3 elliptic genus.
\item Gaberdiel--Hohenegger--Volpato 2010, 2011, 2012 ---
$\Mtf$-symmetry of the K3 sigma model; twined elliptic genera;
Persson--Volpato 2012 transgression classification.
\item Gannon 2016 \emph{Adv.\ Math.}\ 301 --- Positivity of the
$A_n$-dimensions as non-negative virtual $\Mtf$-representations.
\item Gritsenko 1999 \emph{St.\ Petersburg Math.\ J.}\ 11 --- Additive
Jacobi--Siegel lifts at integer weights.
\item Gritsenko--Cl\'ery 2011 \emph{J.\ Reine Angew.\ Math.}\
(arXiv:0812.3962) --- The 8 diagonal-divisor paramodular forms theorem.
\item Gritsenko--Nikulin 1996, 1998 \emph{Invent.\ Math.}\ 116, 132
--- Automorphic forms and BKM superalgebras; rank-3 construction.
\item Igusa 1962 \emph{Amer.\ J.\ Math.}\ 84 --- Ring of Siegel
modular forms, $\Delta_5, \Delta_{10}, \Delta_{12}$ generators.
\item Lorgat 2020 (self, April 2020) --- Borcherds lift of $\phi_{0, 1}$
and GBKM superalgebras for $\Delta_5$; Conjecture 1.
\item Maass 1964 \emph{Nachr.\ Akad.\ Wiss.\ G\"ottingen} --- Explicit
multiplier system of $\Delta_5$.
\item Mukai 1988 \emph{Invent.\ Math.}\ 94 --- Symplectic-automorphism
classification of K3: finite subgroup of $M_{23} \subset \Mtf$.
\item Nikulin 1984 \emph{Proc.\ Steklov Inst.\ Math.}\ 165 --- K3
surfaces with finite automorphism group.
\item Oberdieck 2018 (arXiv:1706.09377) --- Twined reduced DT zeta
function for $(S \times E)/(\ZZ/N\ZZ)$, Conjecture C.
\item Oberdieck--Pandharipande 2016 --- K3$\times$E Igusa cusp form
theorem, proof of $Z^{K3 \times E}_{\mathrm{red}, \DT} = -1/\Phi_{10}$.
\item Oberdieck--Pixton 2018 --- Holomorphic anomaly equations and
Igusa cusp form conjecture, proof reduction.
\item Persson--Volpato 2012 (arXiv:1207.6227) --- $H^3(\Mtf, U(1))$
transgression, anomalous-class identification.
\item Segal 1988 \emph{Oxford Lecture Notes} --- Elliptic cohomology
and modular functors.
\item Moore--Seiberg 1989 \emph{Commun.\ Math.\ Phys.}\ 123 ---
Polynomial equations for RCFT and modular categories; bi-equivariant
torus partition functions.
\item Moore 2004 --- Attractor mechanism, BPS state algebra,
Drinfeld-centre realisation of RCFT torus invariants.
\end{itemize}

\end{document}
